\section{Redes Neuronales Convolucionales}
Una red neuronal convolucional es una arquitectura de red neuronal utilizada
principalmente para tareas de reconocimiento de patrones en estructuras organizadas
en forma de cuadrículas, como es el caso del procesamiento de imágenes y la visión
por computadora.
\medskip

Para ello, hacen uso de una operación matemática conocida como convolución en lugar de
transformaciones lineales densas en al menos una de sus capas. Además, introducen en su
arquitectura conceptos como campos receptivos locales, pesos compartidos y
submuestreo espacial.

\subsection{Convoluciones}
En su forma más general, la convolución es una operación aplicada a dos funciones
definidas sobre una variable independiente. Esta operación expresa el solapamiento
de una función $g$ sobre una función $f$, resultando en una combinación de ambas
mediante la integral del producto punto a punto.
\medskip

La convolución continua entre dos funciones $f(t)$ y $g(t)$ se define como:
\begin{equation}
    (f * g)(t) = \int_{-\infty}^{\infty} f(\tau)\, g(t - \tau)\, d\tau
\end{equation}
donde:
\begin{description}
    \item[$f$] es la señal de entrada.
    \item[$g$] es el núcleo o kernel.
\end{description}
Cabe notar que, en su definición matemática formal, la convolución implica la
inversión del kernel. Sin embargo, en muchas implementaciones prácticas, incluyendo
las redes neuronales convolucionales, se utiliza una operación equivalente a la
correlación cruzada, en la cual no se realiza dicha inversión. A pesar de esta
diferencia, ambas operaciones son funcionalmente similares en este contexto.

Esta operación puede interpretarse como el desplazamiento de una de las funciones
a lo largo del dominio, calculando en cada punto la integral del producto entre
ambas. La convolución mide, por tanto, el grado de similitud entre las funciones
a medida que una se desplaza respecto a la otra.
\medskip

En el caso discreto, esta operación se aproxima mediante sumas finitas, lo cual
permite su aplicación en dominios computacionales como el procesamiento de datos
discretos:

\begin{equation}
    (f * g)(i) = \sum_{\tau} f(\tau)\, g(i - \tau)
\end{equation}

La convolución discreta puede extenderse a dos o más dimensiones, lo que la convierte
en una operación fundamental para el procesamiento de imágenes. Cabe notar que, en su
definición matemática formal, la convolución implica la inversión del kernel. Sin embargo,
muchas implementaciones prácticas, incluyendo las redes neuronales convolucionales, no
realizan dicha inversión, utilizando en su lugar una operación equivalente a la
correlación cruzada.
\medskip

En este contexto, se puede definir una imagen $I$ como la entrada y un kernel
bidimensional $K$, de modo que la operación se expresa como:
\begin{equation}
    (I * K)(i,j) = \sum_{m}\sum_{n} I(i + m, j + n)\, K(m,n)
\end{equation}

Esta operación puede interpretarse como una suma ponderada de la vecindad de $(i,j)$,
donde el kernel $K$ proporciona los pesos, midiendo la similitud entre la región local
de la imagen $I$ y el patrón representado por $K$. La región de la entrada que participa
en este cálculo, definida por el área cubierta por el kernel, corresponde al campo
receptivo local en esa posición. En el contexto del procesamiento de imágenes, este
patrón se conoce como filtro.
\medskip

Una interpretación más intuitiva, en el contexto del procesamiento de imágenes, consiste
en considerar una cuadrícula (el kernel) que se posiciona sobre una región de la imagen
de entrada. En cada posición, se multiplica el valor de cada elemento del kernel con el
valor del píxel correspondiente en la imagen, y la suma de estos productos define el
valor de salida en dicha posición. Este proceso se repite a medida que el kernel se
desplaza a lo largo de toda la imagen.

\begin{figure}[ht]
    \centering
    \includegraphics[width=0.9\textwidth]{diagrams/chapter2/convolution_1.pdf} 
    \caption{Ejemplo de convolución 2D.}
    \label{fig:convolution}
\end{figure}

En el procesamiento de imágenes, los valores asignados a los pesos del kernel pueden utilizarse para
construir filtros que permiten preprocesar los datos, resaltar patrones o extraer características específicas
de la imagen.

\begin{figure}[ht]
    \centering
    \includegraphics[width=1\textwidth]{diagrams/chapter2/convolution_examples.pdf} 
    \caption{Ejemplo de filtros y sus kernels.}
    \label{fig:convolution_example}
\end{figure}


\subsection{Conectividad Local y Compartición de Pesos}

\subsection{Arquitectura de la Red Convolucional}

\subsection{Mapas de Activación y Representaciones}